Quando os elementos estruturais denominados feixes vibram livremente, as suas frequências naturais, $\omega$, dependem da densidade mássica de um feixe, $\rho$, do seu módulo de elasticidade, $E$, e do seu comprimento, $l$, profundidade, $h$ e área da secção transversal, $A$. Se um modelo e um protótipo devem ser construídos com o mesmo material e testados, e seus comprimentos são dimensionados na proporção de $1:5$, como suas frequências naturais se relacionarão? (\textit{Dica:} use a análise dimensional para determinar os vários parâmetros adimensionais que relacionam $\omega$ com as várias propriedades do feixe.)


